% Generated draft: requires mathematical and source review.
\documentclass[11pt]{article}
\usepackage[margin=25mm]{geometry}
\usepackage{fontspec}
\setmainfont{Latin Modern Roman}
\usepackage{amsmath,amssymb,mathtools,unicode-math}
\setmathfont{Latin Modern Math}
\usepackage{enumitem,needspace,xurl,hyperref}
\setlength{\parindent}{0pt}
\setlength{\parskip}{5pt}
\setlength{\emergencystretch}{3em}
\newcommand{\N}{\mathbb{N}}
\newcommand{\Z}{\mathbb{Z}}
\newcommand{\Q}{\mathbb{Q}}
\newcommand{\R}{\mathbb{R}}
\newcommand{\C}{\mathbb{C}}
\newcommand{\F}{\mathbb{F}}
\newcommand{\A}{\mathbb{A}}
\newcommand{\PP}{\mathbb{P}}
\newcommand{\E}{\mathbb{E}}
\newcommand{\eps}{\varepsilon}
\newcommand{\dd}{\,\mathrm{d}}
\begin{document}
\section*{UnsolvedMath ID: 20002058}
\textbf{Source problem code:} catalogue entry
\subsection{Definitions and mathematical statement}
Problem \(34\) asks for other natural tensor-densities satisfying an absent transformation law \((43)\), observing that each gives a conformally invariant natural operator \(L\). The aid reconstructs this as the affine law
\[
 N_{\widehat g}=N_g+L_g\omega,\qquad\widehat g=e^{2\omega}g.
\]
The original sentence ends ``Solutions to''; no further target or classification is supplied. The following is the reported Branson--Gover construction, not an unrestricted completion of that fragment.

In the form-\(Q\) framework, closed \(k\)-forms satisfy, with its weighted-form identifications and normalizations,
\[
 Q_k^{\widehat g}u=Q_k^gu+\delta Q_{k+1}^g d(\omega u).
\]
Let \(\tau\) be a natural closed conformally invariant Pontryagin form of degree \(k\) for which this construction applies. Pairing gives the natural top-form density
\[
 N_g=\langle\tau,Q_k^g\tau\rangle_g\,\dd V_g,
\]
with affine variation
\[
 L_{\tau,g}(\omega)=
 \langle\tau,\delta Q_{k+1}^g d(\omega\tau)\rangle_g\,\dd V_g.
\]
The aid reports that this variation is a formally self-adjoint conformally invariant natural operator, and that \(\int_MN_g\) is conformally invariant on closed manifolds. Closedness of \(\tau\), formal self-adjointness of the form operator and the degree/weight normalization are substantive parts of the construction. An arbitrary chosen closed form would not automatically give a natural density.

The aid's explicit nonzero family is: for \(r\geq1\), take a Riemannian product \(X^{4r}\times Y^{4r+2}\), where \(X\) is Einstein and locally symmetric, \(Y\) has constant curvature, and
\[
 \operatorname{Ric}_X=\lambda_X g_X,\qquad
 \operatorname{Ric}_Y=\lambda_Y g_Y.
\]
Assume the pulled-back degree-\(4r\) Pontryagin character form \(\tau_r\) of \(X\) is nonzero. It is parallel in this setting, and in the aid's standard near-middle normalization,
\[
 Q_{4r}\tau_r=
 \frac{2((4r+2)\lambda_Y-(4r+1)\lambda_X)}{8r+1}\,\tau_r.
\]
In particular \(X=\mathbb{CP}^{2r}\) with its standard locally symmetric Einstein metric and \(Y=S^{4r+2}\) with a round metric give a nontrivial \(Q\)-like affine density in real dimension \(8r+2\), except when
\[
 (4r+2)\lambda_Y=(4r+1)\lambda_X.
\]
The nonzero-form assumption is necessary for a general \(X\); not every locally symmetric factor has the required Pontryagin form. The exceptional scaling ratio concerns this product value and does not classify all possible affine densities. The source's open-ended construction request is answered in the reported framework by examples, not by a uniqueness theorem for every natural tensor-density.
\subsection{Short English statement}
The aid reconstructs the missing law as an affine conformal transformation and supplies Branson--Gover form-\(Q\) constructions paired with natural Pontryagin forms. A product family on \(\mathbb{CP}^{2r}\times S^{4r+2}\) has an explicit eigenvalue and is nonzero except at one scaling ratio. No complete classification is claimed.
\subsection{Sources}
\begin{itemize}
\item Local UnsolvedMath registry, permanent ID \(20002058\), source code \texttt{catalogue entry}; canonical collection \texttt{aim-geometry-notes.json}, notes[395].
\item AIM workshop ``Conformal structure in geometry, analysis, and physics,'' Problem \(34\): \url{https://aimath.org/WWN/confstruct/confstruct.pdf}. Formula \((43)\) and the continuation after ``Solutions to'' are absent. No PDF was opened in this session.
\item The aid attributes the form-\(Q\) construction to Branson and Gover and uses natural Pontryagin character forms. No full authoritative bibliography, endpoint or normalization conventions are supplied; original remarks and literature fields are empty, with no proposer or year.
\item Supplied research attempt \(1\), result \texttt{full\_solution}, run status ``solved,'' rigor label ``proved,'' novelty confidence ``low.'' Registry status ``solved,'' classification \texttt{SOLVED-IN-LITERATURE}, updated 14 August 2026. The claimed explicit contribution is the product eigenvalue family within the known construction.
\end{itemize}
\end{document}
